\section{Introduction}

There are twenty known crystalline phases of ice.
More amorphous forms are known.
In conditions which exist on Earth, only hexagonal ice (ice $I_h$) is relevant.
High pressure phases, however, have been observed in experiment.
Ice-VII is a hydrogen-disordered phase which forms at pressures over 2GPa and temperatures over 270K.
Ice-VIII is the hydrogen-ordered low-temperature (>270K) analogue of Ice-VII.
An interesting recently observed phenomenon is the high-temperature (~480K) phase transition of Ice-VII to Ice-VIIp (plastic Ice-VII).
Molecules in this phase are free to rotate as a liquid while occupying fixed positions.
\par

Model potentials in Molecular Dynamics (MD) simulations are constantly in development. 
Traditional fixed-charge methods are relatively computationally inexpensive, however with improvements of computing resources
modern potentials such as the Atomic Multipole Optimized Energetics for Biomolecular Simulation (AMOEBA) force-field
which explicitly treat polarization effects are now much more common.
The simulation of high-pressure phases of ice with more accurate models should produce valuable results.
\par

+extra background.
\par

\begin{itemize}
    \item{explore ice-vii and ice-viii theory more thoroughly}
    \item{elaborate on the amoeba model}
    \item{previous research/computer simulation of high-press ice}
    \item{relevance of high-press ice}
\end{itemize}
